\section*{Chapitre \ref{ch:g11:prob} --- Probabilités et variables aléatoires}

\begin{solution}{exo:g11:prob:1}
\omterm{def:g11:prob:model}{Modèle uniforme} sur $52$ cartes. Cœurs~: $\frac{13}{52} = \frac14$.
Rois~: $\frac{4}{52} = \frac{1}{13}$. Cœur ou roi~:
$\frac{13 + 4 - 1}{52} = \frac{16}{52} = \frac{4}{13}$ (le roi de cœur
est compté une fois). Ni l'un ni l'autre~: $1 - \frac{4}{13} =
\frac{9}{13}$.
\end{solution}

\begin{solution}{exo:g11:prob:2}
$S = 6$~: les couples $(1,5), (2,4), (3,3), (4,2), (5,1)$, donc
$\P(S = 6) = \frac{5}{36}$.

$S \geq 10$~: trois couples pour $10$, deux pour $11$, un pour $12$~:
$\P(S \geq 10) = \frac{6}{36} = \frac16$.

$S$ pair~: les sommes $2, 4, 6, 8, 10, 12$ ont
$1 + 3 + 5 + 5 + 3 + 1 = 18$ couples~: \omterm{def:g10:proba:distribution}{probabilité}
$\frac{18}{36} = \frac12$.
\end{solution}

\begin{solution}{exo:g11:prob:3}
Les \omterm{def:g10:proba:distribution}{probabilités} somment à $1$~: $p = 1 - 0.3 - 0.2 - 0.4 = 0.1$. Puis
\[
\E(X) = 0.3 \times (-1) + 0.2 \times 0 + 0.4 \times 2 + 0.1 \times 5
= -0.3 + 0.8 + 0.5 = 1 .
\]
\end{solution}

\begin{solution}{exo:g11:prob:4}
$\E(Y) = 2 \times 3 - 5 = 1$~; $\V(Y) = 2^2 \times 4 = 16$~;
$\sigma(Y) = 4$.
\end{solution}

\begin{solution}{exo:g11:prob:5}
Les gains sont $0$, $2$, $5$ euros avec \omterm{def:g10:proba:distribution}{probabilités} $0.5$, $0.3$,
$0.2$~; en soustrayant la mise de $2$ euros, $G$ prend les valeurs
$-2, 0, 3$~:
\begin{center}
\begin{tabular}{c|ccc}
$g$ & $-2$ & $0$ & $3$\\
\hline
$\P(G=g)$ & $0.5$ & $0.3$ & $0.2$
\end{tabular}
\end{center}
$\E(G) = -1 + 0 + 0.6 = -0.4$~: le joueur perd $40$ centimes par partie
en \omterm{def:g11:stat:mean}{moyenne} --- défavorable.
\end{solution}

\begin{solution}{exo:g11:prob:6}
Tirage ordonné~: $\P(X = 2) = \frac35 \times \frac24 = \frac{3}{10}$
(rouge puis rouge)~; $\P(X = 0) = \frac25 \times \frac14 =
\frac{1}{10}$ (bleu puis bleu)~; d'où
$\P(X = 1) = 1 - \frac{3}{10} - \frac{1}{10} = \frac{6}{10}$. Puis
\[
\E(X) = 0 \times \frac1{10} + 1 \times \frac6{10} + 2 \times \frac3{10}
= \frac{12}{10} = 1.2 ,
\]
cohérent avec l'intuition~: $\frac35$ des deux tirages sont rouges en
\omterm{def:g11:stat:mean}{moyenne}.
\end{solution}

\begin{solution}{exo:g11:prob:7}
$\E(G^2) = 0.5 \times 4 + 0.3 \times 0 + 0.2 \times 9 = 3.8$, donc par
la formule de König--Huygens
$\V(G) = 3.8 - (-0.4)^2 = 3.8 - 0.16 = 3.64$ et
$\sigma(G) \approx 1.9$ euros.
\end{solution}

\begin{solution}{exo:g11:prob:8}
Le bénéfice est $80$ moins la réclamation~:
\begin{center}
\begin{tabular}{c|ccc}
$b$ & $80$ & $-920$ & $-4920$\\
\hline
$\P(B=b)$ & $0.94$ & $0.05$ & $0.01$
\end{tabular}
\end{center}
$\E(B) = 75.2 - 46 - 49.2 = -20$~: à ce prix la compagnie \emph{perd}
$20$ euros par police en \omterm{def:g11:stat:mean}{moyenne}. La prime est trop basse pour les
risques couverts (il faudrait qu'elle dépasse $100$ euros pour un
bénéfice espéré positif).
\end{solution}

\begin{solution}{exo:g11:prob:9}
Le score aléatoire est $+1$ avec \omterm{def:g10:proba:distribution}{probabilité} $\frac14$ et $x$ avec
\omterm{def:g10:proba:distribution}{probabilité} $\frac34$~:
\[
\E = \frac14 + \frac{3x}{4} = 0 \iff x = -\frac13 .
\]
Avec la pénalité $-\frac13$, le hasard ne rapporte rien en \omterm{def:g11:stat:mean}{moyenne}.
\end{solution}

\begin{solution}{exo:g11:prob:10}
Deux pièces~: elles coïncident (FF ou PP) avec \omterm{def:g10:proba:distribution}{probabilité}
$\frac24 \times 2 = \frac12$. Le gain du joueur est $+1$ ou $-1$ avec
\omterm{def:g10:proba:distribution}{probabilités} égales~: \omterm{def:g11:prob:expectation}{espérance} $0$ pour les deux joueurs --- le jeu est
équitable.

Trois pièces~: toutes coïncident (FFF ou PPP) avec \omterm{def:g10:proba:distribution}{probabilité}
$\frac{2}{8} = \frac14$. Le joueur reçoit $2$ avec \omterm{def:g10:proba:distribution}{probabilité}
$\frac14$ et paie $1$ avec \omterm{def:g10:proba:distribution}{probabilité} $\frac34$~:
$\E = \frac24 - \frac34 = -\frac14$. Le jeu favorise Anne de $25$
centimes par manche~; il serait équitable si elle payait $3$ euros pour
un triple match.
\end{solution}

\begin{solution}{exo:g11:prob:11}
\emph{1.} L'organisateur reçoit $1000 \times 5 = 5000$ euros et verse
$2000 + 5 \times 200 + 20 \times 25 = 3500$ euros au total~: le profit
est la \emph{constante} $1500$ (pas de hasard une fois tous les billets
vendus). Le gain $G$ d'un seul joueur (lot moins le billet à $5$ euros)~:
\begin{center}
\begin{tabular}{c|cccc}
$g$ & $1995$ & $195$ & $20$ & $-5$\\
\hline\rule{0pt}{11pt}%
$\P(G=g)$ & $\dfrac{1}{1000}$ & $\dfrac{5}{1000}$ & $\dfrac{20}{1000}$
& $\dfrac{974}{1000}$
\end{tabular}
\end{center}

\emph{2.}
\[
\E(G) = \frac{1995 + 5 \times 195 + 20 \times 20 - 974 \times 5}{1000}
= \frac{1995 + 975 + 400 - 4870}{1000} = -1.5 .
\]
Chaque joueur s'attend à perdre $1.50$ euro~; sur $1000$ joueurs c'est
$1500$ euros --- exactement le profit de l'organisateur. Les comptes
s'équilibrent~: ce que les joueurs perdent en \omterm{def:g11:stat:mean}{moyenne} est précisément ce
que l'organisateur encaisse.
\end{solution}

\begin{solution}{pb:g11:prob:1}
\textbf{1.} $G$ prend les valeurs $-2, -1, 0, 1, 2, 3$, chacune
avec la \omterm{def:g10:proba:distribution}{probabilité} $\frac16$~: $\E(G) = \frac{-2 - 1 + 0 + 1 + 2 +
3}{6} = \frac12$ euro.

\textbf{2.} $\E(G^2) = \frac{4 + 1 + 0 + 1 + 4 + 9}{6} =
\frac{19}{6}$~; $V(G) = \frac{19}{6} - \frac14 = \frac{35}{12}
\approx 2.92$~; $\sigma(G) \approx 1.71$.

\textbf{3.} $\E(2G - 1) = 2 \times \frac12 - 1 = 0$~;
$V(2G - 1) = 4\,V(G) = \frac{35}{3}$.

\textbf{4.} $\E(\text{somme}) = 7$~: le prix équitable est de $7$
euros. À $8$ euros, la perte espérée est de $1$ euro par partie.

\textbf{5.} Prime équitable~: $0.05 \times 800 = 40$ euros. En
faisant payer $60$~: bénéfice espéré de $20$ euros par contrat.

\textbf{6.} Partager selon $5:3$ récompense le \emph{passé} ---
alors que l'enjeu appartient à qui gagnera l'\emph{avenir}, et B
peut encore l'emporter. Tout donner à A traite une victoire
probable comme certaine --- la petite chance de B vaut quelque
chose.

\textbf{7.} Manches à jouer~: $3$~; issues (victoire de A ou de B à
chaque manche)~: AAA, AAB, ABA, ABB, BAA, BAB, BBA, BBB. Le joueur
B ne remporte la série qu'en gagnant les trois manches~: la seule
issue BBB. A remporte la série dans $7$ cas sur $8$~: partage
équitable $64 \times \frac78 = 56$ pistoles pour A et $8$ pour B.

\textbf{8.} Jouer les manches supplémentaires, devenues inutiles, ne
change la victoire de personne~: une fois que A a sa sixième
victoire, les manches suivantes ne sont que cérémonie. Mais fixer
l'horizon à exactement $3$ manches rend les $8$ issues
équiprobables --- un dénombrement honnête réclame des histoires de
même longueur, et les rallonger ne coûte rien.

\textbf{9.} À $5$--$5$~: une manche décisive, $32$ pour chacun. À
$5$--$4$~: A obtient $\frac{64 + 32}{2} = 48$. À $5$--$3$~: la
manche suivante mène soit à $64$ (si A la gagne), soit à la
position $5$--$4$ qui vaut $48$~: la part de A est
$\frac{64 + 48}{2} = 56$ --- le $56 : 8$ de Fermat, retrouvé en
remontant depuis la fin.

\textbf{10.} A a besoin de $2$ victoires, B de $3$~: on joue $4$
manches imaginaires ($16$ issues). B remporte la série avec $3$ ou
$4$ victoires~: $4 + 1 = 5$ issues. Partage~:
B~: $64 \times \frac{5}{16} = 20$~; A~: $44$.

\textbf{11.} La valeur présente équitable d'un gain futur incertain
est son \omterm{def:g11:prob:expectation}{espérance} --- la \omterm{def:g11:stat:mean}{moyenne} pondérée par les \omterm{def:g10:proba:distribution}{probabilités} de ce
qui peut advenir. L'assurance (question 5) et la finance
(valorisation de flux incertains) sont cette phrase transformée en
industries~; les casinos aussi, de l'autre côté de la table.

\textbf{12.} L'enjeu finit toujours quelque part~: les deux parts
sont des \omterm{def:g11:prob:rv}{variables aléatoires} de somme constante égale à $64$, donc
par linéarité $\E_A + \E_B = 64$ --- les parts équitables épuisent
l'enjeu, à tout score.

\textbf{13.} Par linéarité~: $\E(X + Y) = 3.5 + 3.5 = 7$, sans
tableau --- et cela reste vrai même pour des dés collés l'un à
l'autre. \omterm{def:g11:prob:variance}{Variance} d'un dé~:
$\E(X^2) - 3.5^2 = \frac{91}{6} - \frac{49}{4} = \frac{35}{12}$~;
pour des dés indépendants,
$V(X + Y) = \frac{35}{12} + \frac{35}{12} = \frac{35}{6}$.

\textbf{14.} Chaque invité fixé récupère son propre chapeau avec la
\omterm{def:g10:proba:distribution}{probabilité} $\frac1n$ (son chapeau est l'un des $n$). En sommant sur
les $n$ invités~: $\E(X) = n \times \frac1n = 1$ --- un invité
chanceux en \omterm{def:g11:stat:mean}{moyenne}, que la soirée compte deux chapeaux ou deux
mille.

\textbf{15.} $n = 2$~: les deux répartitions donnent soit les deux
chapeaux bien rendus ($X = 2$), soit l'échange ($X = 0$)~:
$\E = 1$. $n = 3$~: les six répartitions donnent
$X = 3, 1, 1, 1, 0, 0$~: $\E = \frac{6}{6} = 1$.

\textbf{16.} Le calcul a additionné les \omterm{def:g11:prob:expectation}{espérances} des $n$
indicateurs individuels («~l'invité $i$ a retrouvé son chapeau~»),
et la linéarité additionne les \omterm{def:g11:prob:expectation}{espérances} \emph{sans condition} ---
la dépendance ruinerait un produit ou une \omterm{def:g11:prob:variance}{variance}, jamais une somme
d'\omterm{def:g11:prob:expectation}{espérances}. Voilà le superpouvoir de la linéarité~: aucune
question d'indépendance n'est posée.

\textbf{17.} Les deux gains~: $\E = 0.5 \times 10 - 5 = 0$ et
$0.005 \times 1000 - 5 = 0$. \omterm{def:g11:prob:variance}{Variances} (du gain)~:
A~: $\E(G^2) = 0.5 \times 25 + 0.5 \times 25 = 25$~;
B~: $0.005 \times 995^2 + 0.995 \times 25 \approx 4\,975$~; donc
$\sigma_A = 5$ et $\sigma_B \approx 70.5$. Les gens achètent B~: ce
qu'ils achètent n'est pas l'\omterm{def:g11:prob:expectation}{espérance} (nulle dans les deux cas)
mais la \emph{dispersion} --- une petite chance de traîne droite qui
change une vie.

\textbf{18.} Le premier face au lancer $n$ a la \omterm{def:g10:proba:distribution}{probabilité}
$2^{-n}$ et rapporte $2^n$~: chaque $n$ contribue pour
$2^{-n} \times 2^n = 1$ à l'\omterm{def:g11:prob:expectation}{espérance}, qui vaut donc
$1 + 1 + 1 + \dots$~: elle est infinie. Et pourtant personne de
sensé ne paie $1\,000$ pour jouer une fois~: avec la \omterm{def:g10:proba:distribution}{probabilité}
$\frac{15}{16}$, le gain n'excède pas $16$. L'\omterm{def:g11:prob:expectation}{espérance} résume le
long terme de paris \emph{répétables}~; un jeu unique à traîne
sauvage ne se valorise pas par sa \omterm{def:g11:stat:mean}{moyenne}.

\textbf{19.} \omterm{def:g11:prob:expectation}{Espérances}~: s'assurer, $-20$~; rester à découvert,
$-0.01 \times 1000 = -10$. La stratégie sans assurance gagne en
\omterm{def:g11:stat:mean}{moyenne} --- pour qui peut absorber la perte et répéter le pari de
nombreuses fois. Une famille qui ne survivrait pas au $-1\,000$
n'est pas dans le métier de la \omterm{def:g11:stat:mean}{moyenne}~: elle paie $10$ de plus
pour effacer la \omterm{def:g11:prob:variance}{variance}. Les assureurs moyennent~; les ménages
achètent de la certitude.

\textbf{20.} Née du partage d'un enjeu entre deux joueurs
impatients~; dotée du superpouvoir de la linéarité, qui additionne
les espoirs sans demander la permission à la dépendance~; aveugle au
risque, que la \omterm{def:g11:prob:variance}{variance} --- les loteries, le jeu de
Saint-Pétersbourg, le contrat de la famille --- doit mesurer à sa
place~; et couronnée l'an prochain par la loi des grands nombres,
qui explique pourquoi les \omterm{def:g11:stat:mean}{moyennes} finissent toujours par
encaisser.
\end{solution}
